\documentclass[10pt,a4paper]{amsart}
\usepackage[margin=19mm]{geometry}
\usepackage{amsmath,amssymb,mathtools}
\usepackage[scr=rsfs,cal=euler]{mathalfa}
\usepackage{xfrac}
\usepackage[unicode,hidelinks,bookmarks=false]{hyperref}
\newcommand{\ie}{i.e.\ }
\newcommand{\cf}{cf.\ }
\newcommand{\resp}{respectively\ }
\newcommand{\Subsection}{\subsection}
\providecommand{\nobreakdash}{\nobreak\hbox{-}}
\setlength{\emergencystretch}{2em}
\multlinegap0pt
\usepackage{amssymb}
\usepackage{tikz}
\usepackage{xcolor}
\newtheorem{thm}{Theorem}
\newtheorem{lem}{Lemma}
\newtheorem{claim}{Claim}
\usepackage{cleveref}
\crefname{thm}{Theorem}{Theorems}
\crefname{lem}{Lemma}{Lemmas}
\crefname{claim}{Claim}{Claims}
\newcommand{\co}{\text{co}}
\newcommand{\ola}{\overleftarrow{C}}
\newcommand{\ora}{\overrightarrow{C}}
\title{On the Hamiltonicity, traceability, and toughness of complements of line graphs.}
\author{Adam Mammoliti}
\date{Source: arXiv:2602.00530v2 (CC BY 4.0)}
\begin{document}
\maketitle

\noindent\emph{Source and licence.} On the Hamiltonicity, traceability, and toughness of complements of line graphs.. Version arXiv:2602.00530v2 (CC BY 4.0), distributed by arXiv under the Creative Commons Attribution 4.0 International licence, http://creativecommons.org/licenses/by/4.0/, as recorded in the arXiv licence metadata for this exact version and shown on the abstract page https://arxiv.org/abs/2602.00530v2. Full licence notice: \texttt{licenses/arxiv-2602.00530-CC-BY-4.0.txt}.\par\smallskip

\noindent\textit{Editorial scope.} Counted unit: the article's thm t:main (source0.tex lines 151--153). The article's complete original proof of it is delivered with it (source0.tex lines 641--765, one \texttt{proof} environment of 125 lines, which the article places away from the statement). No mathematics is written, repaired or re-derived: the statement and the proof are the article's own lines, in the article's own notation, and no retained line is substituted or re-flowed.  Retained uncounted as the source-local context the proof needs: lines 459--470; lines 517--524; lines 560--565; lines 580--584; lines 604--606; lines 624--626.  Nothing was omitted from any retained span, so the package carries no omission record.  No retained line contains a citation, so the wrapper carries no bibliography and no bibliography entry.  No retained line contains a figure environment, an image-inclusion command or drawing code, so no image is delivered and the package declares no asset dependency.  Editorial formatting only, in addition: an AMS \texttt{amsart} preamble that restates the article's own declarations and macros used by the retained lines, namely the environments \texttt{thm}, \texttt{lem}, \texttt{claim} on separate counters, together with the standard packages those lines need.  The retained environments print in this document as Theorem 1, Lemma 1, Claim 1 in order of appearance, whereas the article prints the same items with the numbering of its own full document.  The complete native source of the article as distributed in the arXiv e-print for this exact version is bundled unchanged as \texttt{sources/193726/source0.tex}.
\begin{lem}\label{l:disconColineGraph}
If $G$ is a graph and $L=\co(G)$ is disconnected then $G$ is either
\begin{enumerate}
\item[(i)] $K_{1,\lambda}$ and $c(L)=\lambda$ and $\rho(G)=2\lambda$, for some $\lambda \geq 1$; 
\item[(ii)] of type-$I$ and $c(L)=2$ and $\rho(G)=|E(G)|+2$; 
\item[(iii)] $C_4$ and $c(L)=2$ and $\rho(G)=6$;
%\item[(iv)] $K_3$ and $c(L)=3$
\item[(iv)] $F_k$ and $c(L)=3$ and $\rho(G) = k+4$ for some, $k\geq 2$;  %$K_3$ a special case of this
\item[(v)] $K_4^-$ and $c(L) = 3$ and $\rho(G)=8$; 
\item[(vi)] $K_4$ and $c(L)=3$ and $\rho(G)=9$.
\end{enumerate}
\end{lem}

\begin{lem}\label{l:propsOfLongestCycle}
Let $L$ be a graph, $C$ be a longest cycle in $L$ and $H$ be a connected component of $L-V(C)$. Then the following hold.
\begin{itemize}
\item[(i)] If $x \in N_L(H)$ then $x^+,x^- \notin N_L(H)$.
\item[(ii)] If $x,y \in N_L(H)$, then there is no $x^+y^+$-path or  $x^-y^-$-path with all internal vertices in $L-V(C)$.
\item[(iii)] $N_L(H)^+\cup\{x\}$ and $N_L(H)^-\cup\{x\}$ are independent sets in $L$ for all $x \in V(H)$.
\end{itemize}
\end{lem}

\begin{lem}\label{l:common_arg}
For any $x_i\prec_C x_j \preceq_C x_k$ (with $x_i\neq x_k$), %either all of $x_i^+ x_j$, $x_i^- x_j$, $x_i^+ x_k$ and $x_i^- x_k$ are not edges in $L$ 
either
$\{x_i^+ x_j, x_i^- x_j, x_i^+ x_k, x_i^-x_k\}\cap E(L) = \emptyset$ or $x_j^- x_k^+ \notin E(L)$. 
In particular, either  $\{x_i^+ x_j$, $x_i^- x_j\}\cap E(L) = \emptyset$ or $x_j^- x_j^+ \notin E(L)$ for all distinct $i,j \in [d]$.
\end{lem} 

\begin{lem}\label{l:nonXYedgesExtendCycle}
%Let $G$ be a graph and $C$ be a longest cycle in $G$.
%Let $x \in G-C$ and %$x_1 ,\ldots , x_d$ be its neighbours in order in $C$. 
Suppose that $ww^+ \in x_i\overrightarrow{C}x_j$ with $i\neq j$. Then at least one of $wz$ and $w^+z'$ is not an edge in $L$ where $\{z,z'\}= \{x_i^-,x_j^+\}$.
\end{lem}

\begin{lem}\label{l:Hs}
The graphs $H_1,H_2$ and $H_3$ all have a coline graph that is tough but is not Hamiltonian.
\end{lem}

\begin{lem}\label{l:eachCompTrivial}
Let $G$ be a graph whose coline graph $L$ is tough and $C$ be a longest cycle of $L$. Then every connected component of $L-V(C)$ is trivial.
\end{lem}

\begin{thm}\label{t:main}
Let $G$ be a graph and $L$ be its coline graph. If $L$ is tough then $L$ is Hamiltonian unless $G$ is isomorphic to $K_5$, $H_1$, $H_2$ or $H_3$.
\end{thm}

\begin{proof}[Proof of \cref{t:main}]
By \cref{l:Hs} and the fact that $\co(K_5)$ is the Petersen graph, we can assume that $G$ is neither $K_5$, $H_1,H_2$ or $H_3$.
Let $L = \co(G)$ and  $C$ be a longest cycle in $L$ with some fixed orientation. Suppose for a contradiction that $C$ is not Hamiltonian.
By \cref{l:eachCompTrivial}, each connected component of $L-V(C)$ is trivial. 
So let $x$ be an isolated vertex in $L-V(C)$, with degree $d \geq 2$ in $L$, since $L$ is tough and so is 2-connected. Then every neighbour of $x$ in $L$ must be in $C$. So let $N_L(x) = \{x_1,x_2,\ldots, x_d\}$ be the set of neighbours of $x$ in $C$, in order with respect to the orientation of $C$, where 
we take the subscripts of the members of $N_L(x)$ modulo $d$.
Let $v$ be the common endpoint of $x_1^+$ and $x$ and $v'$ be the other endpoint of $x$.

We often interpret the conclusions of Lemmas~\ref{l:propsOfLongestCycle}--\ref{l:nonXYedgesExtendCycle} in terms of edges in $G$, rather than in terms of the vertices in $L$.
%We will directly conclude from Lemmas~\ref{l:propsOfLongestCycle}--\ref{l:nonXYedgesExtendCycle} the natural result for edges in $G$
We will also make constant use of the following facts throughout the proof. 
By \cref{l:propsOfLongestCycle}(iii), $N_L(x)^+\cup \{x\}$ and $N_L(x)^-\cup \{x\}$ are independent sets in $L$, so $x$ is adjacent to the edges in $N_L(x)^+$ and $N_L(x)^-$
Any vertex $y$ in $L-V(C)$ other than $x$ must also be isolated, by \cref{l:eachCompTrivial}. 
Furthermore $y$ cannot be adjacent to two vertices in $N_L(x)^+$ or two vertices in $ N_L(x)^-$, as otherwise $L$ contains either the path $x_i^+ y x_j^+$ or the path $x_i^- y x_j^-$ for some distinct $i, j \in [d]$, contrary to \cref{l:propsOfLongestCycle}(ii).
Also if $uu^+ \in E(C)$ with $u,u^+ \notin N_L(x)$, then the edges $u$ and $u^+$ must be adjacent to $x$ but not each other. Therefore the edges $u$ and $u^+$ must be incident to different endpoints of $x$.

First we prove the following claim that simplifies the remainder of the proof.
\begin{claim}\label{cl:xi+xi-}
$x_i^+x_i^- \notin E(L)$ for all $i \in [d]$.
\end{claim}
\begin{proof}
Suppose for a contradiction and without loss of generality that $x_1^+x_1^- \in E(L)$. Then $x_1^+$ is incident to $v$ and $x_1^-$ is incident to $v'$. Since $N_L(x)^-$ is an independent set in $L$, $x_1^+ \neq x_{2}^-$ and therefore $x_1^{+2} \neq x_{2}$ and so $x_1^{+2}$ is incident to $v'$.
By \cref{l:nonXYedgesExtendCycle} with $(w,w')=(x_1^+,x_1^{+2})$ and 
$(z,z')=(x_1^-,x_j^+)$, $x_j^+$ is adjacent to $x_1^{+2}$ for all $j \neq 1$.
Also, by \cref{l:common_arg}$(j,1,1)$, $x_1$ is adjacent to both $x_j^+$ and $x_j^-$ for all $j \neq 1$. 
Therefore, given that $N_L(x)^+$ is an independent set in $L$ and $x_1$ and $x_1^+$ are not adjacent, $x_j^+$ is adjacent to $x$, $x_1$, $x_1^+$ and $x_1^{+2}$, which is possible only if $x_j^+$ is incident to $v$ and is adjacent to both $x_1$ and $x_1^{+2}$ on the same endpoint for all $j\neq 1$. Clearly only one such edge $x_j^+$ can exist. So $d=2$.
Furthermore
%$x_2$
%Thus, because $x_1^{+2}$ is incident to $v'$, $x_1,x_j^+$ and $x_{1}^{+2}$ are all incident to the same vertex.
%It follows that 
$x_{2}^+$ is not adjacent to $x_1^-$, which implies that $x_2^{+}\neq x_{1}^-$.

If $x_1^{-2} \neq x_{2}^+$, then $x_1^{-2}$ is incident to $v$ but is not $x_{2}^+$.
Then $x_1^{-2}$ is not adjacent to $x_1^{+2}$ and so 
\[
C' = xx_1 x_1^+ x_1^- x_{2}^+ \ora x_1^{-2} x_1^{+2} \ora x_2 x 
\]
is a cycle longer than $C$. So $x_1^{-2} = x_{2}^+$.
If $x_1^{+2} \neq x_2^-$, then $x_2^-$ must be adjacent to $x_1$ and $x_1^-$, which is possible only if $x_2^-$ is incident to $v'$. As $x_1^{+2} \neq x_2^-$, $x_2^-$ cannot be adjacent to $x_2^+$. Hence, given that $x_1^{+2}$ cannot be adjacent to $x_2$, 
\[
C' = x x_1 x_1^+ x_1^- x_2^+ x_2^- \ola x_1^{+2} x_2 x
\]
is a cycle longer than $C$. Consequently $ x_1^{+2} = x_2^-$.
Finally as $x_2$ cannot be adjacent to $x_2^-$ or $x_2^+$, $G$ must be either $H_1$, $H_2$ or $H_3$.
\end{proof}
By \cref{l:propsOfLongestCycle}(iii), $N_L(x)^+\cup \{x\}$ and $N_L(x)^-\cup \{x\}$ are each an independent set in $L$. Therefore as edges in $G$, each of $N_L(x)^+ \cup \{x\}$ and $N_L(x)^- \cup \{x\}$ must form either a star or a triangle, where the latter is possible only if $d=2$.
We separate the proof into cases based on which of these occurs, in each case yielding a contradiction.

\noindent \textbf{Case 1}: The vertices of $N_L(x)^+\cup\{x\}$ and $N_L(x)^-\cup \{ x\}$ correspond to stars on the same vertex. %and either $d\neq 2$ or $N_L(x)^+\neq N_L(x)^-$.

%Let $v$ be the common vertex of both stars and $v'$ the other endpoint of $x$.
Exactly half the vertices  of $x_i \ora x_{i+1}^-$, as edges in $G$, are incident to $v$ for each $i \in [d]$, since the edges corresponding to the vertices in $x_i^+ \ora x_{i+1}^-$ alternate between being incident to $v$ and being incident to $v'$ and the first and final edges are incident to $v$. 
As every vertex in $C$ is in exactly one $x_i \ora x_{i+1}^-$ for some $i\in [d]$, it follows that exactly half of the edges corresponding to vertices in $C$ are incident to $v$.

We show that whenever there is an isolated vertex $y \neq x$ in $V(L)-V(C)$, the edge $y$ must be incident to $v$. Suppose otherwise that $y$ is incident to $v'$. 
The edge $y$ cannot be adjacent to $x_i^+$ for some $i$, as then $y$ is not adjacent to both $x_i^-$ and $x_i$, contradicting \cref{l:propsOfLongestCycle}. On the other hand
$y$ must be adjacent to all but at most 1 edge in $N_L(x)^+$ by \cref{l:propsOfLongestCycle}. Since $N_L(x)^+$ has at least two elements, we have a contradiction. Consequently, $y$ is incident to $v$.

Thus, there are 
\[
\frac{|V(C)|}{2}+|V(L)-V(C)|  = \frac{|V(L)|+(|V(L)|-|V(C)|)}{2}\geq \frac{|V(L)|+1}{2} = \frac{|E(G)|+1}{2}
\] 
edges in $G$ incident to $v$, since $|E(G)|=|V(L)|$ and $|V(L)|-|V(C)| \geq 1$.
However by \cref{l:disconColineGraph}(i), $L$ cannot be tough, since $G$ has a $K_{1,\Delta}$ subgraph and 
\[
|E(G)| < 2\left(\frac{|E(G)|+1}{2} \right)\leq 2\deg_G(v) \leq 2\Delta = \rho(K_{1,\Delta})
\]
edges, a contradiction.

$ $
\\
\noindent \textbf{Case 2}: $d\geq 3$ and the vertices of $N_L(x)^+ \cup \{x\}$ and $N_L(x)^- \cup \{x\}$ correspond to stars but on different endpoints of $x$.

Since $N_L(x)^+ \cup \{x\}$ and $N_L(x)^- \cup \{x\}$ correspond to stars on different endpoints of $x$, we must have that $N_L(x)^+ \cap N_L(x)^- =\emptyset$. 
By \cref{cl:xi+xi-}, the edge $x_j^+$ is adjacent to $x_j^-$ for all $j \in [d]$. Furthermore, because $x_j^+$ is adjacent to at most one edge in $N_L(x)^-$, we have that $x_{i+1}^- x_{j}^+ \in E(L)$ for all $i \in [d]$ and $j\neq i+1$.
Therefore \cref{l:common_arg}$(i,i+1,j)$ for $j \in [d]-\{i,i+1\}$ implies that the edges $x_i^+$ and $x_i^-$ are adjacent to all the elements in $N_L(x)-\{x_i\}$ 
%in $L$. It follows that the edges $x_i^+$ and $x_i^-$ are both adjacent all of the elements in $N_L(x)-\{x_{i}\}$ 
for all $i\in [d]$.
Thus the edge $x_i$ is adjacent to every element in $N_L(x)^\pm-\{x_i^+,x_i^-\}$  for all $i\in [d]$. Since $x_j^+$ and $x_j^-$ are incident to a unique endpoint (that isn't an endpoint of $x$) for each $j$, $x_i$ can be adjacent to at most two pairs $x_j^+$ and $x_j^-$. Therefore, $d \leq 3$ and consequently $d=3$. Furthermore, the elements of $N_L(x)^\pm\cup \{x \}$, as edges, form a $K_5$.
As $G$ cannot be $K_5$, there is either another vertex in $V(L)-V(C)$ other than $x$ or a vertex in $C$ that isn't in  $N_L(x)^\pm $.
Any isolated vertex $y \neq x$ of $L-V(C)$ must, as an edge, be adjacent to $x$ and two elements from each of $N_L(x)^+$ and $N_L(x)^-$, by \cref{l:propsOfLongestCycle}(ii). However, this 
% and isn't adjacent to $x$. So as an edge in $G$, $y$ must be adjacent to $x$ along with all edges except at most one in each of $N_L(x)^+$ and $N_L(x)^-$, which 
is impossible for an edge that is not $x$.
So there is a vertex in $C$ that is not in $N_L(x)^\pm$, say without loss of generality $x_1^{+2} \neq x_2^-$. Then the edge $x_1^{+2}$ is not adjacent to $x_2^+$ and so $x_1^{+2}x_2^+ \in E(L)$. Therefore as $x_1^+ x_3^{-} \in E(L)$ and $x_1^+ x_1^{+2} \in x_3\ora x_2$, \cref{l:nonXYedgesExtendCycle} applied to $(w,w')=(x_1^+,x_1^{+2})$ and $(z,z')=(x_2^+,x_3^-)$, yields a contradiction.
$ $\\

\noindent \textbf{Case 3}: $d=2$ and the edges
$N_L(x)^+ \cup \{x\}$ and $N_L(x)^- \cup \{x\}$ don't both form stars on the same vertex.
%As $d=2$, $N_L(x)^+ \cup \{x\}$ and $N_L(x)^- \cup \{x\}$

Since $N_L(x)^+ \cup \{x\}$ and $N_L(x)^- \cup \{x\}$ are independent sets in $L$ and $d=2$, the elements of $N_L(x)^+ \cup \{x\}$ and $N_L(x)^- \cup \{x\}$ as edges form either stars or triangles. By \cref{cl:xi+xi-} the edges $x_1^+$ and $x_1^-$ are adjacent as are the edges $x_2^+$ and $x_2^-$. Given both these constraints and noting that $N_L(x)^+$ and $N_L(x)^-$ both have cardinality 2 and their elements are incident to $v$ and $v'$ respectively, while it is possible that $x_1^+=x_2^-$ or $x_2^+=x_1^-$, it must be the case that the elements of $N_L(x)^\pm \cup \{x\}$ as edges in $G$
forms a graph $J$ that is isomorphic to $K_3$ when both $x_1^+=x_2^-$ and $x_2^+=x_1^-$, is isomorphic to $F_3$  when exactly one of $x_1^+=x_2^-$ and $x_2^+=x_1^-$ is true and is isomorphic to $K_4^-$ when $x_1^+\neq x_2^-$ and $x_2^+ \neq x_1^-$.
Without loss of generality, suppose that we don't have both $x_1^+ \neq x_2^-$ and $x_2^+=x_1^-$.
%Let $v$ be the common endpoint of $x_1^+$ and $x$ and $v'$ be the other endpoint of $x$.

First suppose that $x_1^- $ is either $x_2^+$ or $x_2^{+2}$ and $x_2^- $ is either $x_1^+$ or $x_1^{+2}$. 
In particular, $V(C)-\{x_1,x_2\} = N_L(x)^\pm$ and so $(V(C)\cup\{x\})-\{x_1,x_2\}= E(J)$.
%, as edges, form a graph $J$ isomorphic to either $K_3$, $K_{1,3}$, $F_3$ or $K_4^- $.
We show that the graph $J'$ formed from the edges $E(G)-\{x_1,x_2\}$  is isomorphic to either $K_{1,\lambda}$ with $\lambda\geq 3$, $F_k$ with $k \geq 2$ or $K_4^-$ as then $G$ has $J'$ as a subgraph and has  $|E(G)|= |E(J')|+2< \rho(J')$ edges, yet $L$ is tough contrary to \cref{l:disconColineGraph}(i), (iv) or (v).
The edges corresponding to the elements of $V(L)-V(C)$ must either form a triangle or a star and each of them must be adjacent to at least one element of $N_L(x)^+$ and one element of $N_L(x)^-$, by \cref{l:propsOfLongestCycle}(ii). 
So considering all possible graphs $J$ and all possible edges corresponding to vertices in $V(L)-V(C)$,  $J'$ isn't isomorphic to either $K_{1,\lambda}$ with $\lambda\geq 3$, $F_k$ with $k \geq 2$ or $K_4^-$ only if $V(L)-V(C) \neq \{x\}$ and %either $J\cong F_3$ with $\deg_J(v)=3$ and some $y \neq x$ in $V(L)-V(C)$ is incident to $v'$ or 
$J$ is isomorphic to $K_4^-$. Furthermore, 
 $x_1^+$ and $x_1^-$ must be incident to $v$ and $x_2^+$ and $x_2^-$ must be incident to $v'$. 
Then any $y \neq x$ in $V(L)-V(C)$ is adjacent to both $x_i^-$ and $x_i^+$ in $L$ for some $i\in [2]$ and thus
%\[
%C' =
%x x_{3-i}x_{3-i}^+ x_{i}^- y x_i^+ x_{3-i}^- x_i
%\]
%Then 
$C' = x x_i x_{i}^+ y x_i^- x_{3-i}^+ x_{3-i} x$ is a cycle longer than $C$, since $x_{3-i}^-$ is the only vertex  from $C$ that isn't in $C'$, a contradiction.

Now suppose that, without loss of generality, $x_1^-$ is neither $ x_2^+$ or $x_2^{+2}$.
Let $z$ be the edge in $\{x_2^+,x_1^-\}$ incident to $v$, such an edge exists since $x_1^+ x_1^- \notin E(L)$ by \cref{cl:xi+xi-}. For any vertex $u \in x_2^{+2}\ora x_1^{-2}$, the edge $u$ cannot be adjacent to $x_1^+$ unless it is incident to $v$, because $u$ is neither $x_1^-$ or $x_2^+$. Consequently $u x_1^+ \in E(L)$ if $u$ is not incident to $v$. Thus, if $x_1^+ =x_2^-$, then 
\[
C' = 
\begin{cases}
 x x_1 x_1^+ x_2^{+2} \ora x_1^- x_2^+ x_2 x &\text{ if } z =x_2^+\\
 x x_1 x_1^- x_2^{+} \ora x_1^{-2} x_1^+ x_2 x &\text{ if } z =x_1^-
\end{cases}
\]
is a cycle longer than $C$, since the edges $x_2^{+2}$ and $x_1^{-2}$ are not incident to $v$ when $z=x_2^+$ and $z=x_1^{-}$, respectively because the edges in $x_2^{+2}\ora x_1^{-2}$ alternate endpoints of $x$. 
Therefore, $x_1^+ \neq x_2^-$.
Since $x_2^{+2}$ and $x_1^-$ are incident to the same endpoint of $x$, $x_2^{+3}$ cannot be $ x_1^-$. As the edges $x_2^{+2}$ and $x_{2}^{+3}$ are incident to different endpoints of $x$, with $x_2^+$ and $x_{2}^{+3}$ incident to the same endpoint of $x$ and $x_2^{+2},x_2^{+3} \notin N_L(x)^\pm$, either the edges corresponding to $N_L(x)^+\cup \{x\}$ and $N_L(x)^-\cup \{x\}$ form stars and $x_2^{+2} x_{1}^+  , x_2^{+3} x_2^- \in E(L)$ or the edges corresponding to $N_L(x)^+\cup \{x\}$ and $N_L(x)^-\cup \{x\}$ form triangles and $x_2^{+2} x_{2}^-  , x_2^{+3} x_{1}^+ \in E(L)$.
Either case contradicts \cref{l:nonXYedgesExtendCycle} applied to  $(w,w')=(x_2^{+2},x_2^{+3})$ and $\{z,z'\} = \{x_1^+,x_2^-\}$.
\end{proof}

\end{document}
